\origpage{31}{}
For the first axiom, let $(A, B)$ be a separation of $\Gamma^+$ of order $< \theta$. Then $(\sigma(A), \sigma(B))$ has order $< q\theta \le ord(\mathcal{T})$ by (1), and so one of $(\sigma(A), \sigma(B)), (\sigma(B), \sigma(A))$ belongs to $\mathcal{T}$. Hence one of $(A, B), (B, A)$ belongs to $\mathcal{T}'$. This verifies the first axiom.

For the second, let $(A_i, B_i) \in \mathcal{T}'$ ($i = 1, 2, 3$). Then $(\sigma(A_i), \sigma(B_i)) \in \mathcal{T}$ ($i = 1, 2, 3$), and so $\sigma(A_1) \cup \sigma(A_2) \cup \sigma(A_3) \ne G$. But
\begin{equation*}
\sigma(A_1) \cup \sigma(A_2) \cup \sigma(A_3) = \sigma(A_1 \cup A_2 \cup A_3)
\end{equation*}
and $\sigma(\Gamma^+) = G$, and so $A_1 \cup A_2 \cup A_3 \ne \Gamma^+$. This verifies the second axiom.

For the third, it suffices by [3, theorem (2.7)] to show that if $e \in E(\Gamma^+)$ and $K_e$ denotes the graph consisting of $e$ and its ends, then $(\Gamma^+ \setminus e, K_e) \notin \mathcal{T}'$, that is, $(\sigma(\Gamma^+ \setminus e), \sigma(K_e)) \notin \mathcal{T}$. If $e \in E(\Gamma)$, or if $e \in S(\Gamma)$ and $e \nsubseteq F_1^* \cup \dots \cup F_r^*$ then $V(\sigma(\Gamma^+ \setminus e)) = V(G)$, and so $(\sigma(\Gamma^+ \setminus e), \sigma(K_e)) \in \mathcal{T}$, since $\mathcal{T}$ satisfies the third axiom. We assume then that $e \in S(\Gamma)$, and $e \subseteq F_1^* \cup \dots \cup F_r^*$. Then $\sigma(K_e) = \mu(e)$, and so $(\sigma(\Gamma^+ \setminus e), \sigma(K_e)) \notin \mathcal{T}$ by (7.1) since $\mu(e)$ is small. This verifies the third axiom.

Hence $\mathcal{T}'$ is a tangle in $\Gamma^+$ of order $\theta$, as required. $\square$

The main result of this section is the following.

\textbf{7.3} \textit{For every surface $\Sigma$ and all integers $p, r \ge 0$ and $q \ge 1$, there exists $\theta \ge 1$ such that, if a graph $G$ has a presentation in $\Sigma$ of depth $\le q$, and defect $\le r$ which surrounds a tangle of order $\ge \theta$, then $G$ has no vital $p$-linkage.}

\textbf{Proof.} Choose $\lambda$ so that (6.3) holds. Choose $\theta_1$ so that (5.2) holds with $\Sigma, p, \theta$ replaced by $\hat{\Sigma}, 2(p + \lambda), \theta_1$. Choose $\theta_2$ so that (4.6) holds, with $\Sigma, \theta', c, \theta$ replaced by $\hat{\Sigma}, \theta_1, c(\Sigma), \theta_2$. Let $\theta = q\theta_2$. We claim that $\theta$ satisfies (7.3). For let $\Gamma, F_1^*, \dots, F_r^*, H_1, \dots, H_r, \mu$ be a presentation of $G$ in $\Sigma$, of depth $\le q$ (we may assume its defect is exactly $r$), surrounding a tangle of order $\ge \theta$. Suppose that $L$ is a vital $p$-linkage in $G$. Let $\Gamma'$ be the union of all paths in $L \cap \Gamma$ with both ends in $bd(\Sigma)$ and no internal vertex in $bd(\Sigma)$ (including one-vertex paths). Then $\Gamma', F_1^*, \dots, F_r^*, H_1, \dots, H_r, \mu$ is a presentation of $G' = \Gamma' \cup H_1 \cup \dots \cup H_r$ in $\Sigma$, of depth $\le q$. Moreover, $L \cap G'$ is a vital linkage in $G'$, by (2.3), and it is a $p$-linkage since each component of $L$ includes at most one component of $L \cap G'$, and it is exhaustive. \fix{We give the details. Let $Z = V(\Gamma) \cap bd(\Sigma)$; then $Z \subseteq V(H_1 \cup \dots \cup H_r)$ by (V1). Let $W$ be the set of vertices $w \in V(\Gamma) \setminus Z$ such that some terminal $t \in V(\Gamma) \setminus Z$ of $L$ is joined to $w$ by a path of $(L \cap \Gamma) \setminus Z$. A vertex of $\Gamma$ not in $Z$ is in no $H_j$, so all its edges in $L$ are in $\Gamma$, by (V1). We claim that $\Gamma' = (L \cap \Gamma) \setminus W$. For let $Q$ be a path of $L \cap \Gamma$ with both ends in $Z$ and no internal vertex in $Z$. Its internal vertices have degree $2$ in $L$, and so $Q \setminus Z$ is a component of $(L \cap \Gamma) \setminus Z$ that contains no terminal of $L$; hence $Q$ has no vertex in $W$. Conversely, let $x \in V(\Gamma) \setminus (Z \cup W)$, and let $P$ be the component of $L$ that contains $x$. Then $x$ is not a terminal of $L$, since $x \notin W$; and following $P$ from $x$ in each direction, every vertex reached before the first vertex in $Z$ is joined to $x$ in $(L \cap \Gamma) \setminus Z$, and so is not a terminal of $L$, and has both its edges of $L$ in $\Gamma$. Hence both directions reach $Z$, and give a path of $L \cap \Gamma$ through $x$, containing both edges of $L$ at $x$, with both ends in $Z$ and no internal vertex in $Z$. An edge of $L \cap \Gamma$ with both ends in $Z$ is itself such a path. This proves our claim. Since $W \subseteq V(\Gamma) \setminus Z$ is disjoint from $H_1 \cup \dots \cup H_r$, it follows that $L \cap G' = L \setminus W$.}

\fix{Let $x, y \notin W$ be vertices of the same component $P$ of $L$, and suppose that some vertex $z$ of $P$ between $x$ and $y$ is in $W$. Then $z$ is joined by a path $R$ of $(L \cap \Gamma) \setminus Z$ to a terminal $t$ of $L$; so $R$ is the subpath of $P$ between $z$ and the end $t$ of $P$, and it passes through $x$ or $y$, which is then in $W$, a contradiction. Hence the subpath of $P$ between $x$ and $y$ is a path of $L \setminus W$; that is, each component of $L$ includes at most one component of $L \setminus W$. So the terminals of $L \setminus W$ in a component $P$ of $L$ lie in one path of $L \setminus W$, and there are at most two of them, and at most one if $|V(P)| = 1$; and $P$ has two terminals of $L$ if $|V(P)| \ge 2$. Hence $L \cap G'$ has at most $p$ terminals. For exhaustiveness, $\Gamma' \subseteq L$; and let $x$ be a terminal of $L \setminus W$ in no $H_j$. Then $x \in V(\Gamma) \setminus (Z \cup W)$, and so $x$ is not a terminal of $L$, and $x$ is incident with two edges of $L$, both in $\Gamma$. If the other end $y$ of one of them were in $W$, then $y \notin Z$, and the edge would join $y$ to $x$ in $(L \cap \Gamma) \setminus Z$, giving $x \in W$. So both edges are in $L \setminus W$, and $x$ is not a terminal of $L \setminus W$, a contradiction.}\fixnote{fix:7.3-linkage-proof-added}{The paper gives only the reason for these two facts. The argument above is the one used in the formalization.} From (6.3), there are at most $\lambda$ $\Sigma$-jumps in $\Gamma'$. Let $X$ be the set of all $v \in V(\Gamma) \cap bd(\Sigma)$ such that $\{v\}$ is the vertex set of a component of $L \cap \Gamma$. Every other component of $L \cap \Gamma$ either \del{is}\fix{includes} a $\Sigma$-jump in $\Gamma'$ \fix{(if both its ends are in $bd(\Sigma)$, its subpath from one end to the next vertex in $bd(\Sigma)$ is such a jump)}, or has an end some vertex $v \in V(\Gamma) \setminus bd(\Sigma)$ which is a terminal of $L$, by (2.2). \fix{Distinct components include distinct $\Sigma$-jumps, since they are disjoint.}\fixnote{fix:7.3-proof}{As printed, the step misses a case: a component of $L \cap \Gamma$ with both ends in $bd(\Sigma)$ may have an internal vertex in $bd(\Sigma)$, and then it is not a $\Sigma$-jump. For example, a path of $L$ may run $a, x, b, y, c$ inside $\Gamma$, with $a, b, c \in bd(\Sigma)$ and $x, y \notin bd(\Sigma)$, and go on in $H_1$ beyond $a$ and $c$. The count below only needs one jump in each such component.} Consequently, $L \cap (\Gamma \setminus X)$ has at most $\lambda + p$ components, and so is a vital $2(\lambda + p)$-linkage in $\Gamma \setminus X$. From (5.2), $\Gamma \setminus X$ has no tangle of order $\ge \theta_1$. Now $\Gamma \setminus X$ is obtained from $\Gamma^+$ by deleting some
